\chapter{Implementation}
\ChapterLabel{Implementation}
\section{Chapter Responsibility}
\label{Section:ImplementationResponsibility}
Make the architecture reproducible through concrete implementation guidance.

\noindent\textbf{Expected result:} A repeatable path from design to a working implementation.

\section{How to Use This Chapter}
Document the selected technology baseline, build and run instructions, coding conventions that enforce architecture, configuration, and definition of done. Keep design rationale in Architecture Overview.
\section{Step 11 Approved Software Choices}
\label{Section:ApprovedSoftwareChoices}
Jack Marino approved Python 3, Tkinter, JSON Lines and the basic controls on
November 23, 2026. This records an implementation direction; it does not claim
application code exists, platform support is demonstrated, or any test passed.
All unresolved values below are TBD, owned by Jack Marino, decision date TBD.
Hardware/device details remain TBD, OMS remains Proposed under
\ArchitectureRef{ARC-002}, and the separate human review remains
\StatusPending\ under \hyperref[Section:HumanReviewChecklist]{its checklist}.

\paragraph{Language and desktop interface.}
Use Python 3 for the simulator and core software, and Tkinter for the desktop
controls \cite{PythonGeneralFAQ,TkinterDocumentation}. Exact Python/Tcl/Tk
versions, packaging, dependency availability and supported OS/CPU/driver
matrix remain TBD. Prefer existing standard-library capabilities where
practical; no external device library or FT245R mode is selected. The approved
controls are authoritative in \hyperref[Table:ApprovedControls]{the UI register}.
GUI refresh scheduling must remain separate from sample/event timestamps;
language or GUI selection does not demonstrate the 10 ms physical-path target.
Threading, queue budgets, scheduling and uncertainty treatment remain TBD.

\paragraph{Recording format.}
Use UTF-8 JSON Lines files with extension \texttt{.jsonl}: one complete JSON
record per line \cite{JSONLinesSpecification}. The approved structure includes
a session header describing format version, acquisition-session identity,
timestamp meaning and whether the original source is simulated or real.
Sample records preserve code, validity, original timestamp, sequence,
acquisition-session identity and timing quality. Represent gaps or recording
errors explicitly where the approved policy requires them. Exact field names,
types, schema/version rules, time units/event/epoch, unknown-time representation,
failure records, partial-file handling and compatibility checks remain TBD.
This is a format selection, not closure of all schema decisions.

During replay, retain the recorded source identity and original sample
metadata; report replay as a separate output mode. A simulated recording must
not become real-hardware evidence when reopened, and replay cannot refresh live
sensor status. Selection controls and basic playback buttons are approved;
pacing, storage/retention limits and failure/loss policies remain TBD.

\paragraph{Implementation and execution boundary.}
Application files, exact launch/build commands and environment checks will be
recorded when implementation exists. No language runtime, GUI package, driver
or dependency is installed by this documentation update. The next planned step
is a usable simulator and software workflow using these approved choices,
with controlled synthetic input. Record software-test evidence with that scope;
real ADC timing, capture, disconnect recovery and platform/driver behavior
remain Pending until demonstrated under their actual test conditions.

\section{Retained Template Implementation Example}
\section{Example}
The implementation uses [language], [framework], and [database]. Build with \texttt{[command]}; run with \texttt{[command]}. Configuration comes from [approved source]; credentials are never stored in source control. Completion requires passing tests, reviewed traceability, and approved evidence.
